\section{Stronger odd-field transfer}\label{upd:transfer}

This appendix gives the odd-characteristic argument with its
coefficient quotient, root interface, native service tariffs, and ordinary
input and output accounting.  The non-elementary native services are named
source theorems with written construction proofs in the native
supplement described below.

\begin{proposition}[Odd-field ordinary product at the stronger exponent]
\label{upd:odd-product}
For each fixed odd prime $p$ and fixed integer $e\geq 1$, there is a
deterministic machine with a fixed finite alphabet and a fixed finite number
of sequential tapes which multiplies two ordinary degree-<$n$ polynomials
over $\F_{p^e}$ in a chosen fixed basis and outputs all $2n-1$ ordinary
coefficients, including zero leading or trailing coefficients.  Its bit
complexity is
\[
  \Oh_{p,e}\!\left(n(\log(n+2))^\gamma
                  (\log\log(n+4))^K\right),
  \qquad \gamma=0.996874195,
\]
for a fixed constant $K$.  Consequently it is
\[
  \Oh_{p,e}\!\left(n(\log(n+2))^{319/320}\right).
\]
The field, basis, constants, and finite machine are fixed before $n$ varies.
No growing-field or uniform-in-$e$ bound is asserted.
\end{proposition}

\begin{proof}
We first construct the machine over $\F_p$.  The fixed-extension reduction
is given at the end.  All logarithms are base two.

\subsection{Cyclotomic host and the ordinary coefficient maps}
\label{upd:host}

Fix a sufficiently large $n$; the finitely many smaller values use a fixed
schoolbook branch.  The unconditional cyclotomic construction of Harvey and
van der Hoeven supplies a squarefree integer $N$ with
\[
  2n<N<4n,\qquad p\nmid N,
\]
its prime factor tuple, and principal-root data~\cite{HVH19Cyclo}.  Write
\[
  L_0=\lceil\log_2 N\rceil,\qquad
  a=\lceil\log_2 L_0\rceil,\qquad R_0=2^{a^2}.
\]
Every prime factor $q_i$ of $N$ lies between $L_0^3$ and $R_0$ once $n$ is
large enough.  The target required by the cyclotomic construction exceeds
its fixed $p^2$ threshold; the direct branch handles the remaining inputs.

The public algorithms for local shapes, permutations, and route settings
are fixed first, with a fixed effective polynomial bound on their record
work.  Starting with the conductor $\alpha_0$ supplied with the root, append
unused prime factors of $N$ to obtain a squarefree conductor $\alpha$ for
which
\[
  D=\varphi(\alpha)\geq (R_0+L_0+5)^A
\]
for a fixed sufficiently large $A$.  Each appended factor contributes
$q_i-1>L_0^3/2$, so $O(\log L_0)$ factors suffice.  With
\[
  m=N/\alpha,\qquad
  \theta=1-\frac{D}{\alpha},
\]
the conductor and the remaining host obey
\begin{equation}
\label{upd:host-shape}
  \log\alpha,\log D=O((\log L_0)^4),\quad
  \theta=O((\log L_0)^2/L_0^3),\quad
  mD=N(1-\theta)=\Theta(N).
\end{equation}
The number of prime axes remaining in $m$ is large enough for the native
dock: the original tuple has at least
$(L_0-O(1))/\log_2 R_0$ axes, and deleting the $O((\log L_0)^2)$ conductor
factors still leaves the required
\[
  h=\left\lfloor
       \frac{L_{\rm dock}}
            {40\lceil\log_2(R_0+1)\rceil}
     \right\rfloor,
  \qquad
  L_{\rm dock}=\lceil\log_2(2M)\rceil,
  \qquad
  M=\prod_{q_i\mid m}(q_i+1).
\]
Here $L_{\rm dock}=L_0-O((\log L_0)^4)=\Theta(L_0)$ and
\[
  \frac{M}{m}=\prod_{q_i\mid m}\left(1+\frac1{q_i}\right)
       =1+O(L_0^{-2}),\qquad MD=\Theta(N).
\]
The large coefficient degree $D$ pays the public setting-generation cost;
the separate relation $MD=\Theta(N)$ pays the data volume.

Let
\[
  S=\F_p[U]/(\Phi_\alpha(U)),\qquad
  A_p=\F_p[I]/(I^2+1),\qquad
  B=A_p[U]/(\Phi_\alpha(U)).
\]
These are literal monic quotients.  The ring $S$ has $D$ scalar
coordinates, while $B$ has $D$ lanes over $A_p$, or $2D$ scalar
coordinates over $\F_p$.  Neither quotient is assumed to be a field.

There are two distinct coefficient maps in the construction.  First,
ordinary input coefficients use the new parameter map $Y_{\rm new}\mapsto U$.
Separately, the old principal root is transported from the original host:
the host parameters satisfy $q_0\mid\alpha_0\mid\alpha\mid N$.
Put $J=\alpha/\alpha_0$. Since $\alpha$ is squarefree,
$\gcd(J,\alpha_0)=1$, and the cyclotomic identity
\[
  \Phi_{\alpha_0}(X^J)=\prod_{d\mid J}\Phi_{\alpha_0d}(X)
\]
implies $\Phi_\alpha(X)\mid\Phi_{\alpha_0}(X^J)$.
Consequently $U_0\mapsto U^J$ defines a unital homomorphism
$\varphi$ from the original coefficient quotient to $S$.
If $\rho_0$ has principal order $N/q_0$ in the original quotient, set
\[
  \omega_0=\rho_0^{\alpha_0/q_0},\qquad
  \omega=\varphi(\omega_0)^J,
  \text{where }\varphi:U_0\longmapsto U^J.
\]
Taking the indicated divisor powers gives principal orders
$N/\alpha_0$ and $m=N/\alpha$, respectively. The homomorphism preserves
the vanishing character sums defining principality. Root transport uses
this homomorphism separately from the ordinary-input coefficient map.

The cyclic coefficient algebra of length $N=\alpha m$ is represented by
\[
  \F_p[X]/(X^N-1)\simeq
  \F_p[U,Z]/(U^\alpha-1,Z^m-1),
  \qquad X\longmapsto U^c Z,
  \quad cm\equiv1\pmod\alpha.
\]
For $k=qm+r$, the monomial $X^k$ maps to
$U^{q+cr\bmod\alpha}Z^r$.  The forward map, its inverse, the short
$\alpha$-direction transpose, and row rotations cost
$O_p(N(1+\log\alpha))$ on fixed banks.  This is the ordinary coefficient
map, not the root lift above.

Both coefficient CRT components are retained.  Since $p\nmid\alpha$,
$U^\alpha-1$ is squarefree over $\F_p$, so
\[
  U^\alpha-1=\Phi_\alpha(U)\,\psi(U),\qquad
  \psi(U)=\frac{U^\alpha-1}{\Phi_\alpha(U)},\qquad
  \gcd(\Phi_\alpha,\psi)=1.
\]
Extended Euclid prepares the two CRT idempotents and inverses.  The
$\psi$-component is reconstructed along with the $\Phi_\alpha$-component;
discarding it would lose ordinary input information.

\subsection{The actual quotient selector and its unit witnesses}
\label{upd:selector}

The algebraic selector is proved inside the actual quotient $B$.  The
conductor $\alpha$ is odd and greater than one, so $D=\varphi(\alpha)$ is
even.  Put $Q=p^D$.  Euler's theorem and oddness give
\[
  Q\equiv1\pmod\alpha,\qquad Q\equiv1\pmod4.
\]
In $B$, $U^\alpha=1$ and $I^2=-1$, hence $U^Q=U$ and $I^Q=I$.  Frobenius
to the $D$-th power is a ring homomorphism in characteristic $p$ and fixes
$\F_p$.  Every element has the monic-basis expression
\[
  x=\sum_{j=0}^{D-1}(a_j+b_jI)U^j,
  \qquad a_j,b_j\in\F_p,
\]
so, by applying Frobenius to the generators and coefficients,
\begin{equation}
\label{upd:quotient-frobenius}
  x^Q=x\quad\text{for every }x\in B.
\end{equation}
This uses the quotient generators themselves.  It does not factor $B$ or
identify it with a field.

For each local coefficient $\mu\in B$, define
\[
  f=\mu^{Q-1},\qquad r=1-f.
\]
Equation~\eqref{upd:quotient-frobenius} gives
\begin{equation}
\label{upd:selector-identities}
  f^2-f=\mu^{Q-2}(\mu^Q-\mu)=0,\qquad
  f\mu=\mu^Q=\mu,\qquad r\mu=0.
\end{equation}
Thus $f$ and $r$ are complementary orthogonal idempotents.  In particular,
the nilpotent remainder needed by the generic word has bound one; the
fixed-field argument uses no Hensel or Newton iteration.

Put $h=2^{-1}\in\F_p$ and $q_U=U^2+1$.  The units used below have explicit
inverses in $S$.  Certainly $U^{-1}=U^{\alpha-1}$.  For $U-1$, let
\[
  c_\alpha=\Phi_\alpha(1),\qquad
  H(X)=\frac{\Phi_\alpha(X)-c_\alpha}{X-1}.
\]
Then $(U-1)H(U)=-c_\alpha$.  Since $\alpha$ is squarefree,
$c_\alpha$ is either $1$ or the prime factor of $\alpha$ when $\alpha$ is
a prime power; either value is nonzero modulo $p$.  Hence
$(U-1)^{-1}=-c_\alpha^{-1}H(U)$.  For $j=1,2$, oddness of $\alpha$ gives
\[
  (1+U^j)\sum_{k=0}^{\alpha-1}(-U^j)^k
      =1-(-U^j)^\alpha=2,
\]
so
\[
  (U+1)^{-1}=h\sum_{k<\alpha}(-U)^k,
  \qquad q_U^{-1}=h\sum_{k<\alpha}(-U^2)^k.
\]
These sums are prepared once by geometric doubling; no inverse is taken as
an algebraic oracle.

The shared canonical parameter and all required inverses are
\begin{align*}
  u&=fU+rI,& u^{-1}&=fU^{-1}-rI,\\
  v_-&=f(U-1)^{-1}-r(I+1)h,&
    &\hspace{-1em}(u-1)v_-=1,\\
  v_+&=f(U+1)^{-1}+r(1-I)h,&
    &\hspace{-1em}(u+1)v_+=1,\\
  c&=f q_U\mu^{Q-2}+r,&
  c^{-1}&=f\mu q_U^{-1}+r.
\end{align*}
Orthogonality of $f,r$, together with~\eqref{upd:selector-identities},
proves each displayed inverse identity.  In particular,
\begin{equation}
\label{upd:associated-q}
  u^2+1=fq_U=c\mu,\qquad cc^{-1}=1.
\end{equation}
The unit $c$ is explicit even when $\mu$ or $u^2+1$ is zero or a zero
divisor.  No factorization, zero test, or residue-field conversion is used.

\subsection[The exact three-C shear word]{The exact three-$C$ shear word}
\label{upd:three-c}

Write $i=I$, $A=(1+i)h$, and $B_0=(1-i)h$, and set
\[
  C=\begin{pmatrix}A&B_0\\ B_0&A\end{pmatrix}.
\]
Then $A B_0=h$, $A^2=ih$, $B_0=-iA$, and $2AB_0=1$.  Let
\[
  b=(u-1)v_+,
  \quad \alpha_1=2B_0uv_+,
  \quad \beta_1=A(u^2+1)v_+,
  \quad \delta_1=B_0(u-1).
\]
Direct multiplication, using $(u+1)b=u-1$, gives the literal three-call
core
\begin{equation}
\label{upd:three-c-core}
  C\begin{pmatrix}u&0\\0&1\end{pmatrix}
   C\begin{pmatrix}b&0\\0&1\end{pmatrix}C
    =\begin{pmatrix}\alpha_1&\beta_1\\0&\delta_1\end{pmatrix}.
\end{equation}
For example,
\[
  C\begin{pmatrix}u&0\\0&1\end{pmatrix}C
   =\begin{pmatrix}ih(u-1)&h(u+1)\\
                     h(u+1)&-ih(u-1)\end{pmatrix};
\]
right multiplication by $\operatorname{diag}(b,1)C$ gives
\eqref{upd:three-c-core}.  Thus the identity is a product of three
positive $C$ factors, not a matrix supplied as a premise.

The two nonzero diagonal pivots in~\eqref{upd:three-c-core} have the
following explicit inverse witnesses:
\[
  \alpha_1^{-1}=A(u+1)u^{-1},\qquad
  \delta_1^{-1}=2A v_-.
\]
Indeed $\alpha_1\alpha_1^{-1}=2AB_0(uv_+)(u+1)u^{-1}=1$ and
$\delta_1\delta_1^{-1}=2AB_0(u-1)v_-=1$.  Normalize the left side of
the core by these units.  The resulting upper shear has parameter
\[
  \kappa=\alpha_1^{-1}\beta_1
       =A^2(u^2+1)u^{-1}
       =ih\,c\mu u^{-1}.
\]
Now put
\[
  z=-2iu c^{-1},\qquad z^{-1}=ih u^{-1}c.
\]
Equation~\eqref{upd:associated-q} gives $zz^{-1}=1$ and $z\kappa=\mu$.
Consequently the four diagonal factors
\[
  D_0=\operatorname{diag}(z\alpha_1^{-1},\delta_1^{-1}),\quad
  D_1=\operatorname{diag}(u,1),\quad
  D_2=\operatorname{diag}(b,1),\quad
  D_3=\operatorname{diag}(z^{-1},1)
\]
are all unit diagonals, and
\begin{equation}
\label{upd:three-c-word}
  D_0 C D_1 C D_2 C D_3
       =\begin{pmatrix}1&\mu\\0&1\end{pmatrix}.
\end{equation}
This one word works for every $\mu$, including $0$ and zero divisors.  Its
three calls and four diagonal slots have fixed topology.  The generic
identity is also the content of the arbitrary-coefficient theorem
\texttt{arbitrary\_threeC\_word\_correct} in
\texttt{OddGeneralThreeCWord.lean}; the displayed host formulas supply its
unit witnesses rather than assuming them.

\subsection{Principal-root differences and the local transforms}
\label{upd:roots}

The cyclotomic host supplies a principal root, whose defining character
sums give explicit differences that are units.  For any $z$ in a
commutative ring define
\[
  G_k(z)=\sum_{j=0}^{k-1}z^j,
  \qquad J_k(z)=\sum_{j=0}^{k-1}(k-j)z^j.
\]
Telescoping (or induction using
$J_{k+1}=J_k+G_{k+1}$) gives
\begin{equation}
\label{upd:weighted-geometric}
  (1-z)J_k(z)=k-zG_k(z).
\end{equation}
If $G_k(z)=0$ and $k$ is invertible, then $J_k(z)k^{-1}$ is an explicit
inverse of $1-z$.  For a principal $k$-th root $\eta$, apply this to
$z=\eta^r$ for $0<r<k$.  This argument works in product rings and
requires no cancellation.  Since all the lengths here divide $m$ and
$p\nmid m$, the needed scalar inverses exist.  This identity proves unit
existence for every divisor length of $m$; it does not evaluate a
length-$m$ sum for each record.  The actual local transforms have prime
lengths $q_i\leq R_0$.  For each required local root difference, direct
evaluation of $J_{q_i}$ uses a fixed polynomial number of operations in
$R_0$ over $B$; an extended-Euclid calculation modulo $\Phi_\alpha$ is
another charged method with a fixed polynomial cost in $D$.  These
preparations are made in public setup, so
the total is $n^{o(1)}$.  No length-$m$ sum is run per local record.

For each prime $q_i\mid m$, let
\[
  E_i=\frac{m}{q_i}
       \left(\left(\frac{m}{q_i}\right)^{-1}\bmod q_i\right),
  \qquad \zeta_i=\omega^{E_i}.
\]
These are the coupled CRT idempotents and roots; they are derived from the
one transported root rather than supplied by an independent root oracle.
For $0\leq j<q_i$, put
\[
  H_j=\prod_{a=1}^j(1-\zeta_i^a),\qquad H_0=1,
  \qquad \beta_j=(-1)^j\zeta_i^{j(j-1)/2}.
\]
Equation~\eqref{upd:weighted-geometric} proves each $H_j$ is a unit.  The
ring-valid Newton-basis identity factors the $q_i$-point Fourier matrix
$F_{r,s}=\zeta_i^{rs}$ as
\[
  F=\mathcal A T\operatorname{diag}(\beta_jH_j^{-1})T^{\mathsf T}\mathcal A,
  \qquad \mathcal A=\operatorname{diag}(H_j),
\]
where $T$ is unit lower triangular Toeplitz with defining sequence
$H_j^{-1}$.  For the entrywise proof, define
\[
  P_0(X)=1,\qquad P_k(X)=\prod_{a=0}^{k-1}(X-\zeta_i^a),\qquad
  g_{j,k}=H_jH_k^{-1}H_{j-k}^{-1}\quad(0\leq k\leq j<q_i).
\]
Set $g_{j,k}=0$ outside this range.  The relation
$XP_k=P_{k+1}+\zeta_i^kP_k$ and
\[
  g_{j,k-1}+\zeta_i^kg_{j,k}
    =H_jH_k^{-1}H_{j+1-k}^{-1}
       \bigl((1-\zeta_i^k)+
                  \zeta_i^k(1-\zeta_i^{j+1-k})\bigr)
    =g_{j+1,k}
\]
for $j+1<q_i$, with the two edge coefficients equal to one, prove by
induction the formal polynomial identity
$X^j=\sum_{k=0}^j g_{j,k}P_k(X)$ for every $j<q_i$.  If $r<k$ then
$P_k(\zeta_i^r)=0$.  If $r\geq k$, direct factorization gives
\[
  P_k(\zeta_i^r)
    =\prod_{a<k}\zeta_i^a(\zeta_i^{r-a}-1)
    =\beta_k H_rH_{r-k}^{-1}.
\]
Substitution in the polynomial identity yields
\[
  \zeta_i^{rs}
    =H_rH_s\sum_{k=0}^{\min(r,s)}
       \beta_k H_k^{-1}H_{r-k}^{-1}H_{s-k}^{-1},
\]
which is exactly the $(r,s)$ entry of the displayed factorization.
Only $H_j$ for $j<q_i$ is inverted; no $H_{q_i}$ inverse, polynomial
equality from finite evaluations, or field cancellation is used.  The
character sums give the inverse transform after multiplication by the
genuine scalar $q_i^{-1}$.

For the local formal convolution let $\mathcal R$ be a nonzero commutative
ring with a specified inverse $h$ of $2$; the actual coefficient ring
$B$ has this property.  At length $v\geq256$, choose $b$ to be the least power of
two with $b^2\geq4v$, put $s=b/2$ and $a=\lceil v/s\rceil$, and choose
$t$ to be the least power of two at least $2a-1$.  For $v<256$ use direct
multiplication.  Since $b^2\geq4v$, $a\leq b/2$ and $t\leq b$.
Minimality of $b$ gives $b<4\sqrt v$.  Also
$t<2(2a-1)<4a\leq8v/b+4$, whence $tb<8v+4b<9v$ because
$4b<16\sqrt v\leq v$ for $v\geq256$.
Finally $b<4\sqrt v\leq v^{3/4}$ in this range.

Work in the free monic quotient
$\mathcal C_b=\mathcal R[Y]/(Y^b+1)$ with generator $y$.  Its free basis
$1,y,\ldots,y^{b-1}$ shows $y^j\ne1$ for $0<j<b$; for
$b\leq j<2b$, $y^j=-y^{j-b}\ne1$ by the same basis independence and
$2\in\mathcal R^\times$.  Thus $y$ has exact order $2b$.
Since $t\leq b$ and both are powers of two, $t\mid b$, so
$\omega=y^{2b/t}$ has exact order $t$.  For $0<r<t$, write
$g=\gcd(r,t)$ and $d=t/(2g)$.  Then $r/g$ is odd and
$(\omega^r)^d=y^{b(r/g)}=-1$, so
\[
  (1-\omega^r)\left(h\sum_{j=0}^{d-1}\omega^{rj}\right)=1.
\]
Every nontrivial root difference is therefore an explicit unit, and
$t^{-1}$ is a unit because $t$ is a power of two.  If $t=2^j$, prepare
$t^{-1}=h^j$ and count each of its applications as a scalar gate.  The
radix-two transform and its inverse are valid over $\mathcal C_b$ with
this genuine normalization.  Split the input polynomials into width-$s$
blocks.  Their outer block product has degree at most $2a-2<t$, and
each inner block product has degree at most $2s-2=b-2$.  Thus neither
transform coordinate wraps; overlap-add recovers the exact product in
$\mathcal R[X]$ and costs $O(ab)=O(v)$.

Counting additions, products and scalar operations, the recursive
supplier has one length-$b$ product for each outer transform point.  The
length-$t$ transforms on $t$ elements of $\mathcal C_b$ use
$O(tb\log(t+2))=O(v\log(v+2))$ coefficient-ring operations; powers of
$y$ are signed coordinate rotations, and the inverse normalization uses
the prepared scalar $t^{-1}$.  Folding, padding, and overlap-add cost
$O(v+ab+tb)=O(v)$.  Therefore
\[
  T(v)\leq tT(b)+K v\log(v+2),\qquad tb<9v,\quad b\leq v^{3/4}.
\]
Since $b\leq v^{3/4}$, we have
$\limsup_{v\to\infty}\log(b+2)/\log(v+2)\leq3/4$.  Choose a fixed
$\eta$ with $9(3/4)^{16}<\eta<1$.  Beyond a fixed threshold, the
recursive coefficient
$9(\log(b+2)/\log(v+2))^{16}$ is at most $\eta$.  Choose $C$ large
enough to absorb the $O(v\log(v+2))$ toll into
$(1-\eta)Cv\log^{16}(v+2)$; increasing $C$ also covers the finite
intervening values and the direct range.  Hence
$T(v)=O(v\log^{16}(v+2))$.  The transform depth sums to
$O(\log(v+2))$ along a recursive path: these parameters give
$t<8v/b+4=O(\sqrt v)$, so each transform has depth
$O(\log(v+2))$, while recursive child logarithms shrink by a factor
at most $3/4$.

Fixing the second convolution vector yields a linear DAG over
$\mathcal R$: evaluate the finite arithmetic circuit in
$\mathcal R[\epsilon]/(\epsilon^2)$, replace a variable input by
$\epsilon x_j$ and the fixed kernel by its constant value, and take the
linear component.  This is a ring identity and uses no field
interpolation.  Replace every fanout by a binary copy tree.  The copied
DAG has size $O(v\log^{16}(v+2))$ and depth $O(\log^2(v+2))$.

The off-diagonal block in the lower-unit-Toeplitz recursion has a
rank-three displacement representation.  Each of its three terms uses
two fixed-kernel convolutions, so one composite chunk DAG contains six
convolutions, all intermediate vectors, copies, pads, summands, and
outputs.  If $u$ is the parent interval and $w$ the larger chunk width,
each chunk's source and target lie in different halves of the original
parent.  The fitting rule counts its actual source and target ports and
every such gate, including the intermediate vectors, and yields
$w=\Omega(u/\log^{16}(u+2))$.  There are
$O(\log^{32}(u+2))$ ragged chunk pairs per recursion level.  Signed dirty
replay executes the DAG, adds its outputs to the target, undoes it, then
executes the source-suppressed DAG and subtracts its outputs before
undoing.  The common dirty contribution cancels, the target receives
exactly the intended chunk map, and every borrowed cell returns to its
incoming value.  Fanout copy trees and the four replay traversals give
$O(\log^2(u+2))$ matching layers per chunk; recursion over
$O(\log(u+2))$ levels therefore uses $O(\log^{35}(u+2))$ rounds on the
original parent coordinates.  The fixed finite range where a composite
chunk does not fit uses direct triangular-shear words.  This printer and its setup lists have
polynomial size in $R_0$; the coefficient arithmetic and record width
are charged separately below.
\subsection{Native services and repriced routing}
\label{upd:native}

Let $V$ denote the whole physical bit volume of a call, including complete
prefixes, gaps, suffixes, independent lanes, coefficient bits,
initialization, head movement, and cleanup.  The native
source supplies the following named theorem interfaces:
\begin{align*}
  I_\tau &: \text{opaque equal binary-field interchange},
  &\tau&=1-\frac{859173879734039}{250000000000000000},\\
  P_\gamma &: C^{\otimes k}\text{ over }
      A_P=(\mathbb Z/P\mathbb Z)[I]/(I^2+1),
  &\gamma&=1-\frac{3125805}{10^9}=0.996874195.
\end{align*}
The $I_\tau$ interface permits arbitrary complete spectators and opaque
payload and costs $O(V(u+1)^\tau\operatorname{polylog}(u+2))$ for $u$-bit
address fields.  The $P_\gamma$ interface applies the literal tensor over
every odd modulus $P$, every coefficient width $s$, and every independent
lane count $\ell$ at cost
\[
  O\!\left(2^k\ell s(k+1)^\gamma
                   \operatorname{polylog}(k+2)\right).
\]
Both tariffs include headers, record encodings, movement, all lanes,
complete spectators, returns, and cleanup.  The interchange proof is
\path{thm:mixedbitswap}.  The phase proof is
\path{thm:mixedbitinteger}.  Their construction uses the pinned source
theorem \path{Stage2Machine.Bound.uniform_bound}~\cite{Stage2Machine26}.
The precise external statement and source revision are recorded in the
supplementary dependency map.
The finite-ring construction retains the contributor attributions recorded
in~\cite{Community26}.

Every aggregate relocation in the strong transfer is charged through
$I_\tau$ as follows.  Suppose two moved objects are contiguous complete
numeric ranges of sizes $A$ and $B$, with $AB\leq M$.  Pad each whole
range separately to a binary power; the combined padded volume is less
than $4AB$.  Apply the lossless binary transpose to this complete rectangle
and crop only after its completed permutation.  Its address width is at
most $\lceil\log_2 M\rceil+2$, so each direction costs
\begin{equation}
\label{upd:rectangle-route}
  O\!\left(VL^\tau\operatorname{polylog}(L+2)\right),
  \qquad L=\lceil\log_2(2M)\rceil.
\end{equation}
This route pads a whole rectangle once.  The earlier row-by-row routing
estimate with exponent $\rho=1-3\cdot2^{-38}$ is larger than $\gamma$ and
does not give~\eqref{upd:rectangle-route}.

At the selected host shape, the native radices satisfy
$L_0^3\leq q_i\leq R_0$, and the per-record width is at least the fixed
polynomial required by the public generators, since $D$ was chosen after
that polynomial was frozen.  The native dock uses disjoint data, dirty
auxiliary, and movement fields; its untouched suffix has at least $0.8L$
bits and pays the polynomial address and gauge work.  The source supplies
the dock in \path{lem:nativedock} and the phase in
\texttt{prop:nativephase}.  The shuffle and ragged routing proofs are
\path{lem:nativeshuffle} and \path{lem:raggednative}.
Controlled permutations use \path{prop:nativecontrolled}; coupled CRT
levels use \path{prop:nativeCRT} and \path{lem:CRTphase}.
The wider host condition
$\log R_0=O((\log L)^2)$ follows from these proofs and the stated dock
hypotheses; it is not inferred from the narrower polynomial-radix
statement alone.

The strong native bounds satisfy
\[
  \tau<\gamma,\qquad
  \left(1-\tau\right)-\frac{3125805}{10^9}
   =\frac{77722629734039}{250000000000000000}>0.
\]
The shuffle has $O(\log R_0)$ endpoint blocks, ragged routing has
$O(\log R_0)$ global levels, and balanced CRT routing has $O(\log t)$
depths, where $t$ is the number of axes.  Replacing every aggregate move
by~\eqref{upd:rectangle-route} charges these as fixed secondary powers of
$\log L$.  The native phase uses $P_\gamma$; no integer-multiplier
conclusion from that source is used.

Each odd $q_i$ axis is completed to the even radix $q_i+1$.  Complete
each local edge matching by pairing its remaining positions, including
the existing dummy coordinate; the remaining count is even.  Give every
such added pair the coefficient $\mu=0$.  Equation~\eqref{upd:three-c-word}
then acts as the identity on arbitrary incoming values, with all three
$C$ calls and all four diagonal slots still present.  The generated local
permutation places pairs in canonical adjacent positions; a separately
priced route-in and route-out implements this change of layout.
Shorter axis schedules are padded to the common tensor schedule with
these zero-weight identity words.

In a synchronized tensor round, the product of local words is three full
$P_\gamma$ calls and four product-diagonal slots.  It is a product of
operator words, not an expansion into exponentially many additive terms.
Inside a word, live values can occupy a dummy or the $I$ helper coordinate;
only its completed shear boundary is used for restoration or cropping.
At completed boundaries the zero shears restore every dummy-pattern class,
and the shear over $S$ restores the helper subspace.  The controlled and
CRT routes preserve those same classes at their completed boundaries.
Thus no crop or helper projection is inserted inside a native word.

All local edge matchings use
$O(L_0R_0\log^{35}(R_0+2))$ coefficient records.  The selector exponent
has $O(D\log p)$ bits.  Root differences, selector powers, diagonals,
archives, shape lists, route settings, and erasure are prepared once per
local record, not once per global sample.  Their total work is a fixed
polynomial in $R_0,D,L_0$, hence $n^{o(1)}$.  The total volume of one
native call is
\[
  V=2MD\,\lceil\log_2 p\rceil=\Theta_p(n).
\]
Including the fixed powers of secondary logarithms, all native phase,
route-in, route-out, and CRT work is therefore
\begin{equation}
\label{upd:native-cost}
  O_p\!\left(n(\log n)^\gamma
                 (\log\log(n+4))^{K_1}\right)
\end{equation}
for a fixed $K_1$.

The native BIT precompiler fixes one moment parameter $t\in(0,1)$ for
which its productive and total profiles satisfy $P_t\leq T_t$; it does
not require this inequality for every $t$.  The application uses
$t=1-173/50000$ and $z=1-3459/10^6>t$.  For rational $z=a/b$, its finite
stopping test is the exact integer comparison
$F(N)^b<B(N)^bN^a$.  The induction sets $B(0)=1$, $F(0)=0$, and
$\ell(0)=0$.  Rank-zero complements are identities and are omitted;
positive-rank recipes have positive productive count.  At height $G$,
the empty-target class has at most $J^G$ columns and is admissible because
$m_{\rm flow}J^G\leq m_{\rm flow}L^G\leq n$.  Thus a source exists and
the normalization denominator is positive.  The finite winning word is
selected before the later input widths vary; it is not an input-sized
advice table or a compiler oracle.

The separated base word has dimension $127^{1024}$ in
\texttt{lem:finitebitflow}.  Its all-column proof and strictification feed
\texttt{thm:finitebitcompiler}, which selects the effective finite winner
and its constants $N_0\geq2$, $B_0$, and $F_0$; the enormous instruction
list is not materialized.  Every paid leaf has rank one.  At width $n$,
a rank-one child has width $\lfloor n/N_0\rfloor$; aside from
nonrecursive old-supplier leaves, each fallback child is no wider, so
every recursive child has width at most
$\lfloor n/N_0\rfloor\leq n/2$.  For the rank-one tree, a safe path
allowance is
\[
  D_0=2(\lceil\log_2(E+1)\rceil+1),
\]
where $E$ is the root digit parameter.  This uses the actual half-width
contraction of the compiled word.

Fix the finite owned/route inventory and choose one eligible odd address
prime $q$ after the winning word is frozen.  Put
\[
  L_{\rm cov}=B_0|\operatorname{GL}_{N_0}(\F_q)|,\qquad
  A_w=q^{N_0^2(w-1)},\qquad W_w=L_{\rm cov}A_w,
  \qquad w=\lceil E^{1/250}\rceil.
\]
Choose the least $r_0$ with $q^{2r_0}\geq W_w^{D_0}$, then increase it
by fewer than $w$ digits to $r$ so that $w\mid E-r$.  Because
\[
  \log_q W_w=\log_q L_{\rm cov}+N_0^2(w-1)=O_{N_0,B_0,q}(w),
\]
we have $r=O_{N_0,B_0,q}(w\log(E+2))=o(E)$.  The finitely many $E$
for which $r\geq E$ use a direct branch.  Exchange the two high
$r$-digit fields and place the right high field before the low head field;
these short moves cost $O(Vr)$.  Pad their joined $q^{2r}$ row values
once to
\[
  R'=W_w^{D_0}\left\lceil\frac{q^{2r}}{W_w^{D_0}}\right\rceil
       <2q^{2r}.
\]
The range supplies an outer factor and $D_0$ consecutive $W_w$ cover
slots.  This is a grouped complete numeric range, not entrywise radix
conversion.

A child fixes only a productive role and low matrix type, consuming one
fresh cover slot while retaining all $A_w$ high matrices as complete
earlier spectators.  At depth $d$, retained-high multiplicity $A_w^d$
times unused stock $R'/W_w^d$ is exactly $R'/L_{\rm cov}^d$.  Thus the
child volume is divided by $L_{\rm cov}$, not by $W_w$; there is no
second high-matrix saving.  Siblings reuse restored stock, parent
inactive banks remain below their stack tops, and the finite owned-bank
inventory travels with the active stream until its prescribed clearing
boundary.  Only the completed root endpoint restores the padded row
support for cropping.  This is the retained-stock argument of
\texttt{lem:threecoverstock}, generalized to the winning finite word by
\texttt{prop:delayedbitcover}.

The finite program gives
\[
  C_0=\frac{3459}{10^6},\qquad C_1=\frac{69}{20000},\qquad
  \frac{F_0}{B_0}<N_0^{1-C_0},\qquad d=C_0-C_1=\frac9{10^6}.
\]
Choose $q$ large enough to avoid every finite denominator and rank
witness and to satisfy $q>256N_0^5/d$.  With
$\sigma=1-C_1$, the fresh-cover density and full fallback give
\[
  \mu_\sigma\leq\frac{F_0}{B_0}N_0^{-\sigma}
       \left(1+\frac{64N_0^5}{q}\right)
   <N_0^{-d}(1+d/4)
   <\frac{1+d/4}{1+d/2}<1,
  \qquad \mu_1\leq\mu_\sigma.
\]
The last inequality uses $N_0^d\geq1+d/2$, since $N_0\geq2$ and
$\ln 2>1/2$.  This is the BIT density estimate.  The separate phase
moment theorem \texttt{lem:strongerphasemoment} supplies $P_\gamma$; it
is not used for the BIT density bound.

The ancestor decoder is charged by \path{lem:rootdecoder}.
Its descriptor work is polynomial in $E$
and is absorbed by the complete $q^w$ target fibres, because
$E\leq w^{250}$.  Let
$a_L=123879734039/10^{15}$ and $\tau_L=1-a_L$, so a port costs
$G(w)=(w+1)^{\tau_L}\operatorname{polylog}(w+2)$ per active volume.
The strict profile bound gives the recurrence estimate
$T(n,w)\leq A n+B n^{1-C_1}G(w)$, with the finite generated-adapter
cost charged linearly.  The one-reserve wrapper costs
\[
  O\!\left(V\left(E/w+(E/w)^{1-C_1}G(w)+w\log(E+2)\right)\right).
\]
With $\vartheta=1/250$, its leading saving is
\[
  a_B=(1-\vartheta)C_1+\vartheta a_L
      =\frac{859173879734039}{250000000000000000},
  \qquad \tau=1-a_B.
\]
Exact rational comparison gives $a_B<\vartheta<1-a_B$, so the stopped
term dominates the other root terms.  The general reserve and finite
profile are used by \path{prop:delayedbitcover}.
The old \path{lem:threecoverstock} is the retained-stock template.
Coefficient preparation is charged by \path{lem:rootdecoder}.
\subsection{Diagonal invoices, the complement, and ordinary output}
\label{upd:invoice}

The four diagonal slots in a tensor round are product tables of already
prepared local weights.  If the prefix products have cardinalities
$M_1,\ldots,M_t$, their construction uses
\[
  \sum_{j=1}^t M_j<2M
\]
$B$-products: in reverse axis order these cardinalities form a geometric
sum with ratio at most $1/2$.  Applying the resulting table to all $M$
sample positions is a separate $M$-product scan.  Copies, inverses,
erasure, and head returns are charged with those scans.  An ordinary
$B$-product costs $O_p(D\operatorname{polylog}L_0)$ by carry-free integer
packing, classical integer multiplication, and reduction modulo the
monic polynomial $\Phi_\alpha$~\cite{SS71,BrentZimmermann10}.  Coefficient
slots have $O(\log D)$ bits; public integer-modulus reciprocals use
$O(\log D)$ precision, while reversed monic reciprocals use their actual
quotient precision.  Thus building and applying one product-diagonal
table costs
\[
  O_p(MD\operatorname{polylog}L_0),
\]
not $O(Mt)$ ring products and not $M$ fresh selector computations.  The
number of rounds is a fixed power of $\log R_0$, so all diagonal slots
contribute $O_p(n(\log\log(n+4))^{K_2})$ for a fixed $K_2$.

It remains to account for the $\psi$-component.  Write
$d_\psi=\alpha-D$ and $v=d_\psi m=N\theta$.  Classical Schönhage--Strassen
arithmetic computes this component with inner stride $2d_\psi$ and long
integer-modulus precision $O(L_0)$, so coefficient bands do not overlap
and integer carries do not cross their slots.  The full $\alpha$-row
monic division by $\psi$ uses $D+1$ reciprocal coefficients; a short
degree-$d_\psi$ product reciprocal is not enough.  Its total cost is
\[
  O_p(N\theta L_0^2\operatorname{polylog}L_0)
      =O_p(N\operatorname{polylog}L_0)
\]
by~\eqref{upd:host-shape}.  All full coefficient-CRT row maps have the
same $O_p(N\operatorname{polylog}L_0)$ bound.  This complement is
computed and recombined.

On the $\Phi_\alpha$ component, embed the valid coefficient data into
$B$, run two completed forward transforms, multiply pointwise in $B$, and
run the inverse. The pointwise stage uses at most $M$ ordinary
$B$-products, hence costs
\[
  O_p(MD\operatorname{polylog}L_0)
     =O_p(n(\log\log(n+4))^{K_3})
\]
for a fixed $K_3$, including the associated scans and head returns.
Initially invalid dummy entries are zero and remain zero
at every completed shear/route boundary; their point products are zero.
The inverse uses the genuine $m^{-1}$ on this reached valid subspace; no
inverse assertion is made for arbitrary nonzero dummy fibres.  Project
away the restored helper coordinate, reconstruct both coefficient CRT
components, invert the ordinary top map, crop only at its completed
boundary, and output the $2n-1$ canonical residues.  Since $N>2n-1$, the
zero-filled cyclic product has no aliasing among degrees $0$ through
$2n-2$.  Input scans, output scans, and final cleanup cost $O_p(n)$.

Combining~\eqref{upd:native-cost}, the diagonal invoice, the complement,
the coefficient maps, and the fixed polynomial setup proves the
$e=1$ case of the proposition.  Each named supplier is a completed
fixed-tape machine.  Their finite bank sets are renamed and reused, and
recursive setup stacks are tape data; joins, returns, and cleanup are
included.  No tape is added per axis, prime, depth, or gate.

Finally fix a basis $\{\beta_1,\ldots,\beta_e\}$ of $\F_{p^e}$ over
$\F_p$ and its structure constants
$\beta_i\beta_j=\sum_k c_{ij,k}\beta_k$.  A linear scan separates each
input into its $e$ prime-field component streams.  Run the $e^2$ ordered
component-pair products serially, reusing the same prime-field banks, then
accumulate the fixed combinations $c_{ij,k}$ and merge the $e$ output
streams coefficient by coefficient.  These codecs and fixed linear
combinations cost $O_{p,e}(n)$, and the $e^2$ calls are a constant factor.
All $2n-1$ output coefficients, including zeros, are retained.  This
proves the fixed-extension case.

For the plain exponent, observe that
\[
  \frac{319}{320}-\gamma
     =\frac{805}{10^9}>0.
\]
For each fixed $K$, the factor $(\log\log(n+4))^K$ is absorbed by
$(\log(n+2))^{805/10^9}$ for sufficiently large $n$; changing the
constant covers the remaining finite range.  This yields the stated
$319/320$ bound.
\end{proof}

\begin{remark}[Source and verification boundary]
\label{upd:boundary}
The native interfaces are \texttt{thm:mixedbitswap} and the phase-only
\texttt{thm:mixedbitinteger} in the accompanying native supplement,
with the constructive adapters and proof labels listed above.
\path{ASSETS.json} binds the supplement and its finite premises;
\path{SOURCE_SELECTION.json} records its derivation from the research
source. Finite checks concern their stated finite identities and do not
replace the all-size arguments. Complete Lean verification remains open.
\end{remark}
